\section{Method}\label{sec:method}

\pvas{} runs one reverse-diffusion trajectory and lets it choose its own
measurements. At each of a few acquisition rounds the sampler pauses and
estimates the per-coordinate posterior variance
$\PV \approx \operatorname{diag}\operatorname{Cov}[\xz \mid \xt]$ from
the denoiser itself, at $2q$ extra calls per round
(Sec.~\ref{sec:method_pv}). This map scores the candidate measurements,
and the budget is spent on the highest-scoring atoms mid-trajectory
(Sec.~\ref{sec:method_active}). Sec.~\ref{sec:method_inst} instantiates
the loop for multi-coil MRI and guided super-resolution.
Sec.~\ref{sec:method_when} places the single trajectory on a
cost--robustness spectrum: a law-of-total-variance decomposition
(Proposition~1) states which uncertainty one trajectory can and cannot
see, and a two-trajectory variant (\pvas-duo) restores the missing
cross-mode term at well below ensemble cost.

\subsection{Problem setup and background}\label{sec:method_setup}

Let $\xz \in \mathbb{R}^d$ or $\mathbb{C}^d$ be observed through a fixed
dictionary $\mathcal{A}$ of \emph{measurement atoms}: indivisible groups
of linear functionals acquired together, such as a k-space phase-encode
column across all coils or a $4{\times}4$ pixel block. For an index set
$M \subseteq \mathcal{A}$,
\begin{equation}
  y \;=\; \Amat_M \xz + n, \qquad n \sim \mathcal{N}\!\big(0, \operatorname{diag}(\sigma_i^2)\big),
  \label{eq:forward}
\end{equation}
with $\Amat_M$ the rows selected by $M$. \emph{Active sensing} starts from
an initial set $M_0$ (the ACS region in MRI; a low-resolution image) and
may add at most $L$ atoms \emph{during} reconstruction.

We use pretrained, unconditional diffusion priors in the VP/DDPM
\cite{ho2020ddpm,song2021ddim} and EDM \cite{karras2022edm}
parameterizations. Both are mapped to a single noise scale
($\xtilde = \xt/\sqrt{\abar}$, $\sigt^2 = (1-\abar)/\abar$ for VP;
$\xtilde = \xt$ for EDM), so that $\xtilde = \xz + \sigt\,\epsilon$ and
the network provides the Tweedie posterior mean
$\mut(\xt, t) \approx \mathbb{E}[\xz \mid \xt]$
\cite{robbins1956empirical,efron2011tweedie}. As in DPS, \pigdm{}, and
DDRM/DDNM \cite{chung2023dps,song2023pigdm,kawar2022ddrm,wang2023ddnm},
denoising is interleaved with data consistency. Our operators are masked
(sub)unitary transforms, so we correct $\hat{\xz} = \mut(\xt,t)$
DDNM-style in the range of $\Amat_M$ and continue the sampler from
\begin{equation}
  \hat{\xz} \;\leftarrow\; \hat{\xz} + \tilde{\Amat}_M^{H}\big( \gain \odot (y - \Amat_M \hat{\xz}) \big),
  \label{eq:dc}
\end{equation}
with $\tilde{\Amat}_M^{H}$ a normalized adjoint (mask transpose in pixel
space; SENSE adjoint in MRI). The gain is the scalar \pigdm-style
$\gain_i = \sigt^2/(\sigt^2 + \sigma_i^2)$ \cite{song2023pigdm}
throughout; a per-coordinate, variance-derived gain exists as an
ablation option (Sec.~\ref{sec:method_active}).

\subsection{Posterior variance from a single trajectory}\label{sec:method_pv}

\paragraph{Tweedie's second-moment identity.}
Differentiating the posterior mean with respect to $\xtilde$ yields
\cite{efron2011tweedie,manor2024posterior}
\begin{equation}
  \operatorname{Cov}[\xz \mid \xt] \;=\; \sigt^2 \, \frac{\partial \mut}{\partial \xtilde}
  \;=:\; \sigt^2 J_t ,
  \label{eq:tweedie2}
\end{equation}
so the denoiser Jacobian carries the full posterior covariance at the
current point. Exact diagonal extraction costs $d$ denoiser calls, so we
estimate it stochastically
\cite{hutchinson1990trace,boys2024tmpd,schioppa2026covaware}.

\paragraph{Hutchinson probes with central finite differences.}
Draw $v \sim \mathcal{N}(0, I)$ and form a Jacobian--vector product by a
central finite difference in the $\xtilde$ domain,
\begin{equation}
  J_t v \;\approx\; \frac{\mut\!\big(\xt + \varepsilon \sqrt{\abar}\, v,\, t\big)
  - \mut\!\big(\xt - \varepsilon \sqrt{\abar}\, v,\, t\big)}{2\varepsilon},
  \label{eq:jvp}
\end{equation}
with no autodiff, any black-box denoiser, and $O(\varepsilon^2)$
truncation error (in EDM: perturbation $\varepsilon\sigt v$ on $\xt$,
difference divided by $2\varepsilon\sigt$). The raw estimate
\begin{equation}
  \PV^{\mathrm{raw}} \;=\; \sigt^2 \,\big(J_t v \odot v\big), \qquad
  \mathbb{E}_v\big[\PV^{\mathrm{raw}}\big] = \sigt^2 \operatorname{diag} J_t,
  \label{eq:pvest}
\end{equation}
is unbiased for the diagonal of any fixed $J_t$ (up to truncation). Its
per-coordinate variance is dominated by off-diagonal row energy,
$\sigt^4\big(2J_{ii}^2 + \sum_{j \neq i} J_{ij}^2\big)$, which motivates
$q$-probe averaging and the atom-level aggregation below. We floor
negative coordinates at zero, trading a small positive bias for lower
variance. For k-space statistics we also use a \emph{positive rank-one}
variant: with $v$ drawn white in k-space and $F$ the unitary Fourier
transform,
\begin{equation}
  \PV^{\mathrm{raw}}_i \;=\; \sigt^2\,\big| F (J_t v) \big|_i^2 \,/\, \overline{|v|^2},
  \label{eq:pvk}
\end{equation}
which estimates the row energies of $F J_t$ and is nonnegative by
construction. When $F$ approximately diagonalizes $J_t$ --- the natural
regime for the radially structured spectra of MR images --- those row
energies reduce to the squared diagonal.

\paragraph{Variance reduction: aggregate over atoms, smooth optionally.}
A one-probe estimate is noisy per coordinate, but the score needs atoms,
not coordinates: $\PV$ is summed over each atom's coordinates ($H$
frequencies of a column; the pixels of a block). This averages exactly
where Eq.~\ref{eq:score} needs it. Optionally, the map is first
projected onto a small, domain-appropriate function class $\mathcal{F}$,
\begin{equation}
  \PV \;=\; \Pi_{\mathcal{F}}\big(\PV^{\mathrm{raw}}\big),
  \label{eq:smooth}
\end{equation}
a Gaussian kernel smoother in pixel space, or in k-space the positive
five-parameter family $\exp\!\big(\sum_{j=0}^{4} \beta_j\, r^{2j}\big)$ of
radial even log-polynomials. Projection buys a further
$O(d/\dim\mathcal{F})$ variance reduction at the price of bias toward
$\mathcal{F}$. That bias can be disqualifying: the radial fit discards
all \emph{angular} k-space structure, precisely what column selection
ranks on. We therefore smooth where benign (D2) and rely on aggregation
alone where not (D1).

\paragraph{Cost.}
By default \pvas{} is \emph{selection-only}: $\PV$ is computed on demand
at the $R_a$ acquisition rounds and nowhere else. Each refresh costs
$2q$ extra denoiser calls, so
\begin{equation}
  \nfe \;=\; T + 2q\,R_a .
  \label{eq:nfe}
\end{equation}
The production budgets are modest. D1 uses $T{=}20$, $q{=}2$, $R_a{=}2$:
$28$ \nfe, $1.4\times$ a passive solve. D2 uses $T{=}100$, $q{=}4$,
$R_a{=}3$: $124$ \nfe, $1.24\times$. The per-step gate variant
(Sec.~\ref{sec:method_active}) instead refreshes every $k$ steps and
costs $T + 2q\lceil T/k \rceil$ ($100$ \nfe{} for D1 at $k{=}1$).

\paragraph{Analytic validation.}
For a Gaussian prior $\xz \sim \mathcal{N}(m, \operatorname{diag}(s^2))$ the
MMSE denoiser is linear and $\operatorname{Cov}[\xz\mid\xt] =
s^2\sigt^2/(s^2+\sigt^2)$ in closed form. With 8 probes averaged over 4
replicate draws of $\xt$ and no smoothing, Eq.~\ref{eq:pvest} recovers
this diagonal with mean ratio $0.994$--$1.005$, per-coordinate correlation
$0.86$--$0.97$, and relative MAE ${\approx}0.12$ across
$t \in \{100,\dots,900\}$ of a $1000$-step VP schedule\RLOCAL{}. For
learned priors no closed form exists; Fig.~\ref{fig:calibration}
evaluates calibration against ensemble and oracle references.

\subsection{Active selection under a measurement budget}\label{sec:method_active}

\paragraph{Score.}
The variance map prices candidate atoms. For an unacquired
$a \in \mathcal{A} \setminus M$,
\begin{equation}
  s(a) \;=\; \sum_{i \in a} \PV_i(t)
  \qquad \text{or, noise-normalized,} \quad
  \sum_{i \in a} \PV_i(t)/\sigma_i^2 ,
  \label{eq:score}
\end{equation}
the (noise-weighted) total posterior uncertainty that acquiring $a$
would resolve. Under the diagonal surrogate the exact expected
uncertainty reduction is $\sum_{i \in a} \PV_i^2/(\PV_i + \sigma_i^2)$;
Eq.~\ref{eq:score} is its noiseless limit. This is the criterion prior
training-free systems evaluate with posterior-sample ensembles
\cite{elata2024adasense,nolan2025ads,barba2025dal,iollo2025bedmri,gu2025gsc};
\pvas{} reads it from one trajectory.

\paragraph{Schedule.}
Selection runs in $R_a$ rounds at equispaced steps in a window
$[t_a, t_b]$, with the budget $L$ split evenly across rounds. The window
sits late enough that $\PV$ reflects image content, and early enough
that the new measurements still shape the trajectory. At each round the
top-$L_r$ unacquired atoms by Eq.~\ref{eq:score} join $M$ and their
measurements are appended to $y$ (read from held-out ground truth;
Sec.~\ref{sec:exp}). D2 re-applies the correction of Eq.~\ref{eq:dc}
with the enlarged mask within the same step, to the \emph{raw} Tweedie
mean; in D1 the enlarged mask first acts at the next step's correction.

\paragraph{Budget and \nfe{} accounting.}
We account budgets in \emph{atoms}, not post-hoc acceleration factors,
which hide differing initializations. Every policy starts from the
identical $M_0$ and ends with identical $|M| = |M_0| + L$ (\emph{matched
budget}). Compute is priced in every table through an \nfe{} column
(Eq.~\ref{eq:nfe}); Tab.~\ref{tab:nfe_cost} collects the costs. The full
domain-agnostic loop is stated as Alg.~S1 in the supplementary material.

\paragraph{A faithful ensemble baseline.}
All prior training-free active-sensing systems obtain their acquisition
signal from an ensemble, so we reimplement that signal inside our own
pipeline and run it head-to-head. $N$ parallel particles of the
\emph{same} passive solver (scalar-gain \pigdm{} in D1; the D2 sampler)
share the measurement state. At each acquisition round, unacquired atoms
are scored by the cross-particle variance of the particles' predicted
measurements --- Tweedie means pushed through the forward operator ---
the measurement-space disagreement signal of ADS and AdaSense
\cite{nolan2025ads,elata2024adasense}. The reconstruction is the
particle mean; the cost is $N\,T$ \nfe. A control with random selection
at the same $N$ isolates particle-mean averaging from selection quality.
This yields an ensemble comparison at matched budget, matched prior, and
matched data consistency, rather than numbers transplanted from a
different prior and protocol.

\paragraph{Variance-gated data consistency (ablation option).}
The same map can also replace the scalar gain in Eq.~\ref{eq:dc}.
Treating the denoiser output as a Gaussian surrogate posterior
$\xz \sim \mathcal{N}\big(\hat{\xz}, \operatorname{diag}(\PV)\big)$ in the
measurement basis, conditioning on a measured $y_i$ gives the exact
scalar Kalman update with gain
\begin{equation}
  \gain_i(t) \;=\; \frac{\PV_i(t)}{\PV_i(t) + \sigma_i^2},
  \label{eq:gate}
\end{equation}
refreshed every $k$ steps; the isotropic surrogate $\PV_i \equiv \sigt^2$
recovers the scalar \pigdm{} gain. We claim no novelty here: diagonal
covariance-aware likelihoods of this structure appear throughout the
posterior-covariance literature
\cite{peng2024covariance,boys2024tmpd,hamidi2024cadps,rissanen2025freehunch}.
Empirically the gate is a byproduct, not a contribution: on top of
selection it is \psnrm-neutral in D2 and below the scalar gain in D1
(Tabs.~\ref{tab:ablation_gate} and~\ref{tab:mri_main}), so it is
\emph{off} by default and reported only as an ablation.

\subsection{Instantiations}\label{sec:method_inst}

\paragraph{D1: multi-coil MRI.}
The forward operator is SENSE sampling,
$(\Amat x)_c = M \odot F(S_c \odot x)$, on fastMRI brain data
\cite{zbontar2018fastmri}. Coil sensitivities $S_c$ are estimated from
the always-acquired ACS region by the standard windowed low-resolution
method \cite{uecker2014espirit}. Per-coil noise levels, estimated from
the four signal-free corner patches of measured k-space (minimum over
corners), modulate a global per-frequency noise map, giving
$\sigma_{c,i}^2$ in the gain of Eq.~\ref{eq:dc}. In production that map
is initialized \emph{uniform} at the mean corner variance, consistent
with noise white in k-space, so only the per-coil structure is
exercised; frequency-dependent noise is supported. Atoms are full
k-space columns across all coils. The prior is an unconditional EDM
model \cite{karras2022edm} over complex-valued slices. Probes follow
Eq.~\ref{eq:pvk} with $v$ white in k-space ($q{=}2$, central
differences), computed only at the $R_a{=}2$ acquisition rounds
($28$ \nfe{} at $T{=}20$). Production uses the raw map --- no
normalization, \emph{no radial smoothing}: variance reduction comes from
the per-column sum of the noise-normalized score (Eq.~\ref{eq:score};
under the uniform noise map this ranking coincides with the plain sum),
and the radial fit of Eq.~\ref{eq:smooth} remains available but is off
by default. Rounds sit by default in the window $[0.5, 0.85]$. After
each corrected mean the sampler re-noises,
$\xt \leftarrow \hat{\xz} + \sigma_{t-1}\epsilon$. In-loop JSENSE-style
refinement of sensitivity magnitudes is optional and stated per
experiment (Sec.~\ref{sec:exp}).

\paragraph{D2: guided super-resolution on natural images.}
The low-resolution image --- a box average over $4{\times}4$ blocks ---
is always observed and enforced by exact back-projection of the
box-average residual; the budget buys full-resolution $4{\times}4$
blocks. Acquired blocks are noiseless in the main configuration, so
their consistency is exact replacement; a noisy-measurement variant
appears in the ablations of Sec.~\ref{sec:exp}. The prior is a
pretrained pixel-space DDPM \cite{ho2020ddpm} on CIFAR-10
\cite{krizhevsky2009cifar} sampled with deterministic DDIM
\cite{song2021ddim} ($\eta{=}0$), reusing the model's
$\epsilon$-prediction for the direction term so corrections propagate at
full strength \cite{wang2023ddnm}. The estimator uses the unbiased form
of Eq.~\ref{eq:pvest} with $q{=}4$ probes averaged over channels,
computed only at the $R_a{=}3$ acquisition rounds ($124$ \nfe{} at
$T{=}100$). The map is then Gaussian-smoothed ($\sigma{=}1$\,px) and
block-summed; rounds sit by default in the window $[0.45, 0.75]$.
Per-experiment settings, and CelebA-HQ results
\cite{karras2018progressive}, are in Sec.~\ref{sec:exp} and the
supplement.

Finally, removing the LR observation from D2 leaves pure block inpainting
from a sparse random mask: identical machinery, drastically weaker
conditioning. We use it deliberately as a \emph{boundary-condition probe}
of where single-trajectory variance stops being a valid acquisition
signal (Secs.~\ref{sec:method_when} and \ref{sec:exp}).

\subsection{One, two, or $N$ trajectories}\label{sec:method_when}

The estimator targets $\operatorname{Cov}[\xz \mid \xt]$: the posterior
spread \emph{given the current trajectory state}. An acquisition policy
ideally wants $\operatorname{Cov}[\xz \mid y]$. The gap between the two
is exactly a mode-commitment term.

\smallskip\noindent\textbf{Proposition 1} (total variance across
modes)\textbf{.} \emph{Let $m$ be any mode assignment for the posterior,
i.e., a discrete latent with
$p(\xz \mid y) = \sum_m p(m \mid y)\, p(\xz \mid y, m)$. Then,
coordinate-wise,}
\begin{equation}
  \operatorname{Var}[\xz \mid y]
  \;=\; \underbrace{\mathbb{E}_{m \mid y}\big[\operatorname{Var}[\xz \mid y, m]\big]}_{\text{within-mode}}
  \;+\; \underbrace{\operatorname{Var}_{m \mid y}\big(\mathbb{E}[\xz \mid y, m]\big)}_{\text{cross-mode}} .
  \label{eq:ltv}
\end{equation}
\emph{A single-trajectory Tweedie estimate tracks only the first term.
The disagreement of the Tweedie means of $K$ independent trajectories,
each committed to a mode drawn from $p(m \mid y)$, is an unbiased
estimator of the second term for any $K \ge 2$.}

The proof is the law of total variance, instantiated (four lines in the
supplement); we claim no depth beyond the instantiation. Its reading is
what matters. A reverse trajectory commits to one mode as it denoises,
and its Tweedie variance measures the curvature around that mode --- a
one-sample estimate of the within-mode term \cite{manor2024posterior}.
The cross-mode term is invisible to any single trajectory by
construction. Particle and sample ensembles
\cite{nolan2025ads,elata2024adasense} estimate the sum; that is what
their $N\times$ compute buys.

\paragraph{When is one trajectory enough?}
Strong conditioning. When the always-available measurements (ACS lines
in MRI; the LR preview in D2) pin down the coarse structure,
$p(\xz \mid y)$ is effectively unimodal: the cross-mode term is small,
and the within-mode map carries what selection needs --- only a correct
\emph{ordering} of atom scores, a weaker demand than calibrated
magnitudes. When conditioning is weak, the trajectory commits to one of
several modes inside the acquisition window. Its local curvature can
then rank atoms contrary to the true error: in the from-scratch
inpainting probe, PV--oracle selection overlap falls below chance
(Sec.~\ref{sec:exp}). We report this failure as a boundary, not a
footnote.

\paragraph{\pvas-duo: restoring the cross-mode term at $K{=}2$.}
Proposition~1 also prescribes the minimal fix: estimate the second term
with the smallest possible ensemble. \pvas-duo runs \emph{two}
synchronized trajectories on the shared measurement state and scores
atoms by the sum of both terms,
\begin{equation}
  s_{\mathrm{duo}}(a) \;=\; \sum_{i \in a} \Big( \PV_i + \tfrac{1}{2}\, d_i^{\,2} \Big),
  \qquad d = \hat{\xz}^{(1)} - \hat{\xz}^{(2)},
  \label{eq:duo}
\end{equation}
where $\PV$ comes from probes on the first trajectory (within-mode) and
$d_i^2/2$ is the unbiased two-sample variance of the trajectories'
Tweedie means (cross-mode). The reconstruction is the trajectory mean.
The result is a cost ladder rather than a binary choice: one trajectory
at $1.2$--$1.4\times$ passive \nfe{} when conditioning is strong;
\pvas-duo at $2.2$--$2.5\times$ when it is not; an $N$-particle ensemble
at $N\times$ when even two modes undersample the posterior.
Sec.~\ref{sec:exp} evaluates the ladder under matched budget.
